\section{Wire exponents and the scale in a sum identity}
\label{sec:foundations}

This section records the elementary algebra needed to interpret a
wire identity as an inequality of distances. We follow
\cite[Definitions A.1 and A.4, pp.~19--20]{Anthropic2026}.
A tensor on a finite nonempty set $I$ of parties belongs to
$\bigotimes_{v\in I}\kk^{n_v}$, with each $n_v\ge1$.
Products group indices by party. A restriction has the form
$s=(\bigotimes_v A_v)t$ for linear maps $A_v$.
A cover of $s$ by $t_1,\ldots,t_A$ is an equality
$s=\sum_j s_j$ with each $s_j$ a restriction of $t_j$.
The empty sum is allowed and is a cover of the zero tensor.

\begin{definition}\label{def:price}
A spectral point is a nonnegative function $\phi$ on tensors with
$\phi(1)=1$, $\phi(st)=\phi(s)\phi(t)$, and
$\phi(s)\le\sum_j\phi(t_j)$ for every cover.
For distinct parties $u,v$, the wire $\langle q\rangle_{uv}$
is a rank-$q$ identity matrix, with one-dimensional indices elsewhere.
Its exponent is defined by
$\phi(\langle q\rangle_{uv})=q^{d_\phi(u,v)}$.
A price is a pointwise limit of functions $c\,d_\phi$, with $c\ge0$.
On an infinite set, every pullback to a finite set must be a price.
\end{definition}

\begin{lemma}[Wire exponents]\label{lem:wire}
The exponents exist, belong to $[0,1]$, and, with
$d_\phi(u,u)=0$, define a pseudometric.
For every nonempty $P\subseteq I$, the diagonal tensor
$\langle q\rangle_P=\sum_{i=1}^q\bigotimes_{v\in P}e_i$
has value $q^{g_\phi(P)}$ with $g_\phi(P)\in[0,1]$.
For every tree $T$ on $P$,
$g_\phi(P)\le\sum_{\{u,v\}\in E(T)}d_\phi(u,v)$.
\end{lemma}

\begin{proof}
This is \cite[Lemma A.2, p.~19]{Anthropic2026}; we give its
normalization details. The function $f(q)=\phi(\langle q\rangle_P)$
is multiplicative and nondecreasing on positive integers, by the
product and restriction identities for diagonal tensors.
It is at least $1$, since every nonzero tensor restricts to $1$.
For integers $a,b>1$, set $r_n=\lfloor n\log a/\log b\rfloor$.
The inequalities $b^{r_n}\le a^n<b^{r_n+1}$ give
$r_n\log f(b)\le n\log f(a)\le(r_n+1)\log f(b)$.
Dividing by $n$ and taking the limit shows that
$\log f(a)/\log a$ is independent of $a$. Covering the diagonal
tensor by $q$ rank-one tensors bounds the exponent by $1$.

Contracting the middle indices of
$\langle q\rangle_{uv}\langle q\rangle_{vw}$ by the equality
functional produces $\langle q\rangle_{uw}$.
Restriction monotonicity gives the triangle inequality; symmetry
is immediate. For a tree, each vertex keeps only equal incident
indices and outputs their common value. Connectedness makes all
surviving indices equal, so the resulting tensor is
$\langle q\rangle_P$. This proves the tree bound.
A singleton diagonal tensor is equivalent to the nonzero scalar
tensor, so its exponent is $0$.
\end{proof}

We also write $U_q(P)=\langle q\rangle_P$ for the diagonal, or GHZ,
tensor, and $U_q=U_q(I)$ when the party set is understood.
When a distinction is needed, we call a point from
Definition~\ref{def:price} \emph{relaxed}. An \emph{ordinary normalized
spectral point} is such a point that is additive under local direct
sums and satisfies $\psi(U_q(I))=q$ for every positive integer $q$.
Theorem~\ref{thm:normalization} explains the relation between the two
notions.

For a wire product $T=\prod_e\langle q_e\rangle_{u_ev_e}$ put
$L_d(T)=\sum_e(\log_2q_e)d(u_e,v_e)$. Its load vector $w(T)$
has coordinate $\sum_{e:\{u_e,v_e\}=\{u,v\}}\log_2q_e$
at the unordered pair $\{u,v\}$.
Thus $\log_2\phi(T)=L_{d_\phi}(T)=\langle w(T),d_\phi\rangle$.

\begin{proposition}[Finite-cover inequality]\label{prop:cover}
If a wire product $S$ is covered by $T_1,\ldots,T_A$, with $A\ge1$,
then
\[
 2^{L_{d_\phi}(S)}\le\sum_{j=1}^A2^{L_{d_\phi}(T_j)},\qquad
 L_{d_\phi}(S)\le\max_j L_{d_\phi}(T_j)+\log_2A.
\]
If $c\,d_\phi=e$, with $c>0$, and
$\Delta=L_e(S)-\max_jL_e(T_j)>0$, then
$\Delta\le c\log_2A$.
Writing instead $d_\phi=\lambda e$ gives
$\lambda\Delta\le\log_2A$.
\end{proposition}

\begin{proof}
Use the cover axiom, multiplicativity, and the bound of a sum by
$A$ times its largest summand. Take logarithms and then scale.
\end{proof}

The multiplying scale $c$ is bounded below by $\Delta/\log_2A$,
whereas $\lambda$ in $d_\phi=\lambda e$ is bounded above by
$\log_2A/\Delta$. For $A=1$, positive saving is impossible.
For $A>1$, an additive $\log_2A$ cannot simply be discarded in a
limit defining a price.

\begin{lemma}[Homogeneous families of covers]\label{lem:homogeneous}
Suppose $S_n$ is covered by $T_{n,1},\ldots,T_{n,A_n}$ on a fixed
finite party set. Let $b_n\to\infty$, $\log_2A_n=o(b_n)$, and
$b_n^{-1}w(S_n)\to w$.
If the normalized loads of all $T_{n,j}$ approach, uniformly in
$j$, a fixed finite set $\{v_1,\ldots,v_r\}$, then every price $d$
satisfies
\[
 \langle w,d\rangle\le\max_{\beta\le r}\langle v_\beta,d\rangle.
\]
\end{lemma}

\begin{proof}
Fix a spectral point, divide the finite-cover inequality by $b_n$,
and take the limit. Uniform approximation of the loads justifies
the maximum. The resulting inequality is continuous and homogeneous,
so it passes to nonnegative multiples and pointwise limits.
The limit in $n$ precedes the limit defining the price.
\end{proof}

\begin{example}\label{ex:tensor-power-selector}
The $n$th power of an $A$-term cover has saving $n\Delta$ and at
most $A^n$ summands. The elementary inequality remains
$\lambda\Delta\le\log_2A$. Taking powers alone, without organizing
the summands or introducing another identity, does not satisfy
the subexponential hypothesis of Lemma~\ref{lem:homogeneous}.
\end{example}
